% Source: https://artofproblemsolving.com/community/c3h3108511_joshys_1000th_post
% Stable identity: aops-topic-3108511-q06
We have $\angle A=60^\circ $, $\angle ACM=10^\circ $ and $\angle ABN=20^\circ $. By the sine rule in $ ABN $, $ ACM $ and $ ABC $, \[\dfrac{AN}{AM}=\dfrac{\sin20^\circ\sin50^\circ\sin110^\circ}{\sin100^\circ\sin70^\circ\sin10^\circ}=\dfrac{\sin80^\circ}{\sin40^\circ}.\]For the last identity cancel $\sin110^\circ=\sin70^\circ $ and use $\sin20^\circ=2\sin10^\circ\cos10^\circ $, $\sin50^\circ=\cos40^\circ $ and $ 2\sin40^\circ\cos40^\circ=\sin80^\circ=\cos10^\circ $. A triangle with included angle $ 60^\circ $ and this ratio of its adjacent sides is similar, by SAS, to a triangle with other angles $ 80^\circ,40^\circ $. Thus $\angle AMN=80^\circ $. Also $\angle AMC=110^\circ $, and the two rays lie inside that angle, so $\angle NMC=110^\circ-80^\circ=\boxed{30^\circ}$.
\begin{asy}
usepackage("amsmath");
size(7cm);
pair A=(0,0),B=(1,0),C=(0.6133407985,1.0623374253),M=(0.2266815969,0),Nn=(0.1736481777,0.3007674664);draw(A--B--C--cycle);draw(C--M--Nn--B);label("$A$",A,SW);label("$B$",B,SE);label("$C$",C,N);label("$M$",M,S);label("$N$",Nn,W);
\end{asy}
